\documentclass[11pt]{article}
\usepackage[utf8]{inputenc}
\usepackage[T1]{fontenc}
\usepackage[spanish,mexico]{babel}
\usepackage{amsmath,amssymb}
\usepackage{geometry}
\usepackage{enumitem}
\usepackage{multicol}
\usepackage{fancyhdr}
\geometry{letterpaper,margin=1.7cm}
\setlength{\parindent}{0pt}
\pagestyle{fancy}
\fancyhf{}
\lhead{Álgebra Lineal}
\rhead{Fluidez Aritmética --- Versión C}
\cfoot{\thepage}

\begin{document}
\begin{center}
{\Large\bfseries Prueba Integrada de Fluidez Aritmética para Álgebra}\\[1mm]
{\large Versión C --- Diagnóstico preliminar}
\end{center}

\textbf{Nombre:} \rule{9cm}{0.4pt}\hfill
\textbf{Fecha:} \rule{3cm}{0.4pt}

\medskip
\textbf{Instrucciones.} Resuelva mentalmente cada reactivo y escriba únicamente
el resultado. No utilice calculadora. Simplifique toda fracción hasta obtener
una fracción irreducible. Trabaje con rapidez, pero procure mantener la precisión.

\medskip
\textbf{Diseño.} Esta versión integra los contenidos de las versiones A y B,
pero omite el bloque específico de máximo común divisor. Conserva ejercicios
de cálculo inmediato, divisibilidad, signos, jerarquía, simplificación,
fracciones, decimales, porcentajes y operaciones mixtas.

\medskip
\textbf{Tiempo:} \rule{2.5cm}{0.4pt}\hfill
\textbf{Contestados:} \rule{1.5cm}{0.4pt}\hfill
\textbf{Aciertos:} \rule{1.5cm}{0.4pt}

\subsection*{I. Aritmética inmediata}
\begin{multicols}{2}\begin{enumerate}[leftmargin=*,start=1]
\item $\displaystyle 7+6=$ \rule{2.1cm}{0.4pt}
\item $\displaystyle 15-8=$ \rule{2.1cm}{0.4pt}
\item $\displaystyle 6\cdot7=$ \rule{2.1cm}{0.4pt}
\item $\displaystyle 42\div7=$ \rule{2.1cm}{0.4pt}
\item $\displaystyle 14+17=$ \rule{2.1cm}{0.4pt}
\item $\displaystyle 25-16=$ \rule{2.1cm}{0.4pt}
\item $\displaystyle 9\cdot6=$ \rule{2.1cm}{0.4pt}
\item $\displaystyle 18+7-5=$ \rule{2.1cm}{0.4pt}
\item $\displaystyle 4+7+8=$ \rule{2.1cm}{0.4pt}
\item $\displaystyle 30-8-7=$ \rule{2.1cm}{0.4pt}
\item $\displaystyle 3\cdot5+4=$ \rule{2.1cm}{0.4pt}
\item $\displaystyle 24\div6-2=$ \rule{2.1cm}{0.4pt}
\item $\displaystyle 5+3\cdot4=$ \rule{2.1cm}{0.4pt}
\item $\displaystyle 2(5+3)=$ \rule{2.1cm}{0.4pt}
\item $\displaystyle 18\div(6-3)=$ \rule{2.1cm}{0.4pt}
\end{enumerate}\end{multicols}
\subsection*{II. División mental y divisibilidad por $2$, $3$, $5$ y $7$}
\begin{multicols}{2}\begin{enumerate}[leftmargin=*,start=16]
\item $\displaystyle 18\div3=$ \rule{2.1cm}{0.4pt}
\item $\displaystyle 27\div3=$ \rule{2.1cm}{0.4pt}
\item $\displaystyle 33\div3=$ \rule{2.1cm}{0.4pt}
\item $\displaystyle 42\div3=$ \rule{2.1cm}{0.4pt}
\item $\displaystyle 51\div3=$ \rule{2.1cm}{0.4pt}
\item $\displaystyle 25\div5=$ \rule{2.1cm}{0.4pt}
\item $\displaystyle 35\div5=$ \rule{2.1cm}{0.4pt}
\item $\displaystyle 45\div5=$ \rule{2.1cm}{0.4pt}
\item $\displaystyle 65\div5=$ \rule{2.1cm}{0.4pt}
\item $\displaystyle 21\div7=$ \rule{2.1cm}{0.4pt}
\item $\displaystyle 35\div7=$ \rule{2.1cm}{0.4pt}
\item $\displaystyle 49\div7=$ \rule{2.1cm}{0.4pt}
\item $\displaystyle 63\div7=$ \rule{2.1cm}{0.4pt}
\item $\displaystyle 54\div3=$ \rule{2.1cm}{0.4pt}
\item $\displaystyle 75\div5=$ \rule{2.1cm}{0.4pt}
\item $\displaystyle 77\div7=$ \rule{2.1cm}{0.4pt}
\item $\displaystyle 84\div7=$ \rule{2.1cm}{0.4pt}
\item $\displaystyle 105\div7=$ \rule{2.1cm}{0.4pt}
\item $\displaystyle 48\div2=$ \rule{2.1cm}{0.4pt}
\item $\displaystyle 66\div2=$ \rule{2.1cm}{0.4pt}
\end{enumerate}\end{multicols}
\subsection*{III. Signos, potencias y jerarquía}
\begin{multicols}{2}\begin{enumerate}[leftmargin=*,start=36]
\item $\displaystyle -3+8=$ \rule{2.1cm}{0.4pt}
\item $\displaystyle 5-9=$ \rule{2.1cm}{0.4pt}
\item $\displaystyle -4-3=$ \rule{2.1cm}{0.4pt}
\item $\displaystyle 6-(-3)=$ \rule{2.1cm}{0.4pt}
\item $\displaystyle (-3)(4)=$ \rule{2.1cm}{0.4pt}
\item $\displaystyle (-5)(-2)=$ \rule{2.1cm}{0.4pt}
\item $\displaystyle 24\div(-6)=$ \rule{2.1cm}{0.4pt}
\item $\displaystyle 3^2+4=$ \rule{2.1cm}{0.4pt}
\item $\displaystyle 2^3-3=$ \rule{2.1cm}{0.4pt}
\item $\displaystyle 5^2-20=$ \rule{2.1cm}{0.4pt}
\item $\displaystyle 2^4\div4=$ \rule{2.1cm}{0.4pt}
\item $\displaystyle 3(2^2)-5=$ \rule{2.1cm}{0.4pt}
\item $\displaystyle 4+2(5-3)=$ \rule{2.1cm}{0.4pt}
\item $\displaystyle 3(6-2)-5=$ \rule{2.1cm}{0.4pt}
\item $\displaystyle 2^3+3^2=$ \rule{2.1cm}{0.4pt}
\end{enumerate}\end{multicols}
\subsection*{IV. Simplificación de fracciones}
\begin{multicols}{2}\begin{enumerate}[leftmargin=*,start=51]
\item $\displaystyle \frac{14}{21}=$ \rule{2.1cm}{0.4pt}
\item $\displaystyle \frac{18}{27}=$ \rule{2.1cm}{0.4pt}
\item $\displaystyle \frac{25}{35}=$ \rule{2.1cm}{0.4pt}
\item $\displaystyle \frac{28}{42}=$ \rule{2.1cm}{0.4pt}
\item $\displaystyle \frac{30}{45}=$ \rule{2.1cm}{0.4pt}
\item $\displaystyle \frac{35}{49}=$ \rule{2.1cm}{0.4pt}
\item $\displaystyle \frac{42}{63}=$ \rule{2.1cm}{0.4pt}
\item $\displaystyle \frac{45}{60}=$ \rule{2.1cm}{0.4pt}
\item $\displaystyle \frac{50}{70}=$ \rule{2.1cm}{0.4pt}
\item $\displaystyle \frac{54}{72}=$ \rule{2.1cm}{0.4pt}
\item $\displaystyle \frac{21}{27}=$ \rule{2.1cm}{0.4pt}
\item $\displaystyle \frac{25}{30}=$ \rule{2.1cm}{0.4pt}
\item $\displaystyle \frac{35}{45}=$ \rule{2.1cm}{0.4pt}
\item $\displaystyle \frac{42}{55}=$ \rule{2.1cm}{0.4pt}
\item $\displaystyle \frac{49}{63}=$ \rule{2.1cm}{0.4pt}
\end{enumerate}\end{multicols}
\subsection*{V. Operaciones con fracciones}
\begin{multicols}{2}\begin{enumerate}[leftmargin=*,start=66]
\item $\displaystyle \frac12+\frac13=$ \rule{2.1cm}{0.4pt}
\item $\displaystyle \frac34-\frac12=$ \rule{2.1cm}{0.4pt}
\item $\displaystyle \frac23+\frac16=$ \rule{2.1cm}{0.4pt}
\item $\displaystyle \frac56-\frac13=$ \rule{2.1cm}{0.4pt}
\item $\displaystyle \frac35+\frac1{10}=$ \rule{2.1cm}{0.4pt}
\item $\displaystyle \frac23+\frac14=$ \rule{2.1cm}{0.4pt}
\item $\displaystyle \frac56-\frac14=$ \rule{2.1cm}{0.4pt}
\item $\displaystyle \frac23\cdot\frac35=$ \rule{2.1cm}{0.4pt}
\item $\displaystyle \frac57\cdot\frac{14}{15}=$ \rule{2.1cm}{0.4pt}
\item $\displaystyle \frac34\div\frac12=$ \rule{2.1cm}{0.4pt}
\item $\displaystyle \frac67\div\frac9{14}=$ \rule{2.1cm}{0.4pt}
\item $\displaystyle \frac{14}{15}\cdot\frac5{21}=$ \rule{2.1cm}{0.4pt}
\item $\displaystyle \frac12+\frac13+\frac16=$ \rule{2.1cm}{0.4pt}
\item $\displaystyle 2-\frac12-\frac14=$ \rule{2.1cm}{0.4pt}
\item $\displaystyle \frac23+\frac12-\frac16=$ \rule{2.1cm}{0.4pt}
\item $\displaystyle \frac12\div\frac14=$ \rule{2.1cm}{0.4pt}
\item $\displaystyle \frac34\div\frac38=$ \rule{2.1cm}{0.4pt}
\item $\displaystyle \left(\frac34-\frac14\right)\div\frac12=$ \rule{2.1cm}{0.4pt}
\end{enumerate}\end{multicols}
\subsection*{VI. Decimales, porcentajes y operaciones mixtas}
\begin{multicols}{2}\begin{enumerate}[leftmargin=*,start=84]
\item $\displaystyle 0.5+0.25=$ \rule{2.1cm}{0.4pt}
\item $\displaystyle 1.2+0.8=$ \rule{2.1cm}{0.4pt}
\item $\displaystyle 2.5-1.5=$ \rule{2.1cm}{0.4pt}
\item $\displaystyle 0.4(5)=$ \rule{2.1cm}{0.4pt}
\item $\displaystyle 0.5(18)=$ \rule{2.1cm}{0.4pt}
\item $\displaystyle 25\%\text{ de }20=$ \rule{2.1cm}{0.4pt}
\item $\displaystyle 50\%\text{ de }18=$ \rule{2.1cm}{0.4pt}
\item $\displaystyle 10\%\text{ de }70=$ \rule{2.1cm}{0.4pt}
\item $\displaystyle 20\%\text{ de }30=$ \rule{2.1cm}{0.4pt}
\item $\displaystyle 75\%\text{ de }8=$ \rule{2.1cm}{0.4pt}
\item $\displaystyle \frac{18}{24}+\frac14=$ \rule{2.1cm}{0.4pt}
\item $\displaystyle \frac{21}{35}+\frac25=$ \rule{2.1cm}{0.4pt}
\item $\displaystyle \frac{28}{42}-\frac16=$ \rule{2.1cm}{0.4pt}
\item $\displaystyle \frac35+\frac{15}{25}=$ \rule{2.1cm}{0.4pt}
\item $\displaystyle \frac47\cdot\frac{21}{8}=$ \rule{2.1cm}{0.4pt}
\item $\displaystyle \frac56\div\frac{10}{9}=$ \rule{2.1cm}{0.4pt}
\item $\displaystyle 1-\left(\frac23-\frac16\right)=$ \rule{2.1cm}{0.4pt}
\end{enumerate}\end{multicols}

\newpage
\section*{Clave de respuestas --- uso del profesor}
\subsection*{I. Aritmética inmediata}
\begin{multicols}{4}\begin{enumerate}[leftmargin=*,start=1]
\item $\displaystyle 13$
\item $\displaystyle 7$
\item $\displaystyle 42$
\item $\displaystyle 6$
\item $\displaystyle 31$
\item $\displaystyle 9$
\item $\displaystyle 54$
\item $\displaystyle 20$
\item $\displaystyle 19$
\item $\displaystyle 15$
\item $\displaystyle 19$
\item $\displaystyle 2$
\item $\displaystyle 17$
\item $\displaystyle 16$
\item $\displaystyle 6$
\end{enumerate}\end{multicols}
\subsection*{II. División mental y divisibilidad por $2$, $3$, $5$ y $7$}
\begin{multicols}{4}\begin{enumerate}[leftmargin=*,start=16]
\item $\displaystyle 6$
\item $\displaystyle 9$
\item $\displaystyle 11$
\item $\displaystyle 14$
\item $\displaystyle 17$
\item $\displaystyle 5$
\item $\displaystyle 7$
\item $\displaystyle 9$
\item $\displaystyle 13$
\item $\displaystyle 3$
\item $\displaystyle 5$
\item $\displaystyle 7$
\item $\displaystyle 9$
\item $\displaystyle 18$
\item $\displaystyle 15$
\item $\displaystyle 11$
\item $\displaystyle 12$
\item $\displaystyle 15$
\item $\displaystyle 24$
\item $\displaystyle 33$
\end{enumerate}\end{multicols}
\subsection*{III. Signos, potencias y jerarquía}
\begin{multicols}{4}\begin{enumerate}[leftmargin=*,start=36]
\item $\displaystyle 5$
\item $\displaystyle -4$
\item $\displaystyle -7$
\item $\displaystyle 9$
\item $\displaystyle -12$
\item $\displaystyle 10$
\item $\displaystyle -4$
\item $\displaystyle 13$
\item $\displaystyle 5$
\item $\displaystyle 5$
\item $\displaystyle 4$
\item $\displaystyle 7$
\item $\displaystyle 8$
\item $\displaystyle 7$
\item $\displaystyle 17$
\end{enumerate}\end{multicols}
\subsection*{IV. Simplificación de fracciones}
\begin{multicols}{4}\begin{enumerate}[leftmargin=*,start=51]
\item $\displaystyle \frac23$
\item $\displaystyle \frac23$
\item $\displaystyle \frac57$
\item $\displaystyle \frac23$
\item $\displaystyle \frac23$
\item $\displaystyle \frac57$
\item $\displaystyle \frac23$
\item $\displaystyle \frac34$
\item $\displaystyle \frac57$
\item $\displaystyle \frac34$
\item $\displaystyle \frac79$
\item $\displaystyle \frac56$
\item $\displaystyle \frac79$
\item $\displaystyle \frac{42}{55}$
\item $\displaystyle \frac79$
\end{enumerate}\end{multicols}
\subsection*{V. Operaciones con fracciones}
\begin{multicols}{4}\begin{enumerate}[leftmargin=*,start=66]
\item $\displaystyle \frac56$
\item $\displaystyle \frac14$
\item $\displaystyle \frac56$
\item $\displaystyle \frac12$
\item $\displaystyle \frac7{10}$
\item $\displaystyle \frac{11}{12}$
\item $\displaystyle \frac7{12}$
\item $\displaystyle \frac25$
\item $\displaystyle \frac23$
\item $\displaystyle \frac32$
\item $\displaystyle \frac43$
\item $\displaystyle \frac29$
\item $\displaystyle 1$
\item $\displaystyle \frac54$
\item $\displaystyle 1$
\item $\displaystyle 2$
\item $\displaystyle 2$
\item $\displaystyle 1$
\end{enumerate}\end{multicols}
\subsection*{VI. Decimales, porcentajes y operaciones mixtas}
\begin{multicols}{4}\begin{enumerate}[leftmargin=*,start=84]
\item $\displaystyle 0.75$
\item $\displaystyle 2$
\item $\displaystyle 1$
\item $\displaystyle 2$
\item $\displaystyle 9$
\item $\displaystyle 5$
\item $\displaystyle 9$
\item $\displaystyle 7$
\item $\displaystyle 6$
\item $\displaystyle 6$
\item $\displaystyle 1$
\item $\displaystyle 1$
\item $\displaystyle \frac12$
\item $\displaystyle \frac65$
\item $\displaystyle \frac32$
\item $\displaystyle \frac34$
\item $\displaystyle \frac12$
\end{enumerate}\end{multicols}
\end{document}
